% =========================================================
% Vector Spaces --- Main File
% =========================================================

\begin{remark*}[Toolkit --- Vector Spaces]
\textbf{Concept family:} The abstract vector space axioms and their immediate
consequences.

\medskip
\textbf{Atomic items in this section:}
\begin{itemize}
  \item Definition: Vector Space (10 axioms)
  \item Definition: Vector and Point
  \item Definition: Real Vector Space
  \item Definition: Complex Vector Space
  \item Theorem: Unique Additive Identity
  \item Theorem: Unique Additive Inverse
  \item Theorem: $0v = 0$
  \item Theorem: $a0 = 0$
  \item Theorem: $(-1)v = -v$
\end{itemize}

\medskip
\textbf{Key predicate:}
\[
  \operatorname{VectorSpace}(V, \mathbf{F}, +, \cdot).
\]
The 10 axioms are closure, commutativity, associativity, additive identity,
additive inverse, scalar closure, both distributive laws, scalar associativity,
and the scalar identity.
\end{remark*}

% ---------------------------------------------------------

\begin{definitionbox}{Definition (Vector Space)}
\begin{definition}[Vector Space]
\label{def:vector-space}
A \textbf{vector space} over $\mathbf{F}$ is a set $V$ together with an
addition on $V$ and a scalar multiplication
\[
\mathbf{F} \times V \to V
\]
such that the following properties hold for all $u,v,w \in V$ and all
$a,b \in \mathbf{F}$:
\begin{enumerate}[label=(\roman*)]
  \item $u + v \in V$;
  \item $u + v = v + u$;
  \item $(u + v) + w = u + (v + w)$;
  \item there exists $0 \in V$ such that $v + 0 = v$;
  \item for every $v \in V$ there exists $-v \in V$ such that
        $v + (-v) = 0$;
  \item $av \in V$;
  \item $a(u+v) = au + av$;
  \item $(a+b)v = av + bv$;
  \item $a(bv) = (ab)v$;
  \item $1v = v$.
\end{enumerate}
\end{definition}
\end{definitionbox}

\begin{remark*}[Standard quantified statement]
\[
\begin{aligned}
&V \text{ is a vector space over } \mathbf{F} \\
&\quad\Longleftrightarrow\quad
\text{$V$ is closed under addition and scalar multiplication,} \\
&\qquad\text{addition on $V$ is commutative and associative,} \\
&\qquad\text{$V$ has an additive identity and additive inverses,} \\
&\qquad\text{and scalar multiplication satisfies the two distributive laws,} \\
&\qquad\text{the associative law for scalar multiplication, and the identity law.}
\end{aligned}
\]
\end{remark*}

\begin{remark*}[Definition predicate reading]
\[
\operatorname{VectorSpace}(V,\mathbf{F},+,\cdot)
\]
\end{remark*}

\begin{remark*}[Negated quantified statement]
\[
\neg\,\operatorname{VectorSpace}(V,\mathbf{F},+,\cdot)
\;\Longleftrightarrow\;
\text{at least one of the ten axioms fails for some } u,v,w \in V,\; a,b \in \mathbf{F}.
\]
\end{remark*}

\begin{remark*}[Interpretation]
A vector space is a set whose elements can be added together and multiplied
by scalars, with these operations satisfying the natural algebraic rules
familiar from $\mathbf{F}^n$.
\end{remark*}

\begin{dependencies}
\begin{itemize}
  \item \hyperref[def:fn]{Definition: $\mathbf{F}^n$} (motivating example)
  \item \hyperref[def:addition-in-fn]{Definition: Addition in $\mathbf{F}^n$}
  \item \hyperref[def:scalar-multiplication-in-fn]{Definition: Scalar Multiplication in $\mathbf{F}^n$}
\end{itemize}
\end{dependencies}

% ---------------------------------------------------------

\begin{definitionbox}{Definition (Vector and Point)}
\begin{definition}[Vector and Point]
\label{def:vector-and-point}
An element of a vector space $V$ is called a \textbf{vector}.

When $V = \mathbf{F}^n$, an element of $\mathbf{F}^n$ may also be called
a \textbf{point}.
\end{definition}
\end{definitionbox}

\begin{remark*}[Standard quantified statement]
\[
\begin{aligned}
&v \in V \;\Longrightarrow\; v \text{ is a vector}, \\
&x \in \mathbf{F}^n \;\Longrightarrow\; x \text{ is a point of } \mathbf{F}^n.
\end{aligned}
\]
\end{remark*}

\begin{remark*}[Interpretation]
The word \emph{vector} emphasizes the algebraic role of an element of $V$.
In $\mathbf{F}^n$, the same object can also be viewed geometrically as a point
with coordinates in $\mathbf{F}$.
\end{remark*}

\begin{dependencies}
\begin{itemize}
  \item \hyperref[def:vector-space]{Definition: Vector Space}
  \item \hyperref[def:reals]{The Real Numbers}
\end{itemize}
\end{dependencies}

% ---------------------------------------------------------

\begin{definitionbox}{Definition (Real Vector Space)}
\begin{definition}[Real Vector Space]
\label{def:real-vector-space}
A \textbf{real vector space} is a vector space over $\mathbb{R}$.
\end{definition}
\end{definitionbox}

\begin{remark*}[Standard quantified statement]
\[
V \text{ is a real vector space}
\;\Longleftrightarrow\;
V \text{ is a vector space over } \mathbb{R}.
\]
\end{remark*}

\begin{remark*}[Definition predicate reading]
\[
\operatorname{RealVectorSpace}(V)
\;\colonequiv\;
\operatorname{VectorSpace}(V,\mathbb{R},+,\cdot).
\]
\end{remark*}

\begin{remark*}[Interpretation]
In a real vector space, the scalars used in scalar multiplication are real numbers.
\end{remark*}

\begin{dependencies}
\begin{itemize}
  \item \hyperref[def:vector-space]{Definition: Vector Space}
  \item \hyperref[def:reals]{The Real Numbers}
\end{itemize}
\end{dependencies}

% ---------------------------------------------------------

\begin{definitionbox}{Definition (Complex Vector Space)}
\begin{definition}[Complex Vector Space]
\label{def:complex-vector-space}
A \textbf{complex vector space} is a vector space over $\mathbb{C}$.
\end{definition}
\end{definitionbox}

\begin{remark*}[Standard quantified statement]
\[
V \text{ is a complex vector space}
\;\Longleftrightarrow\;
V \text{ is a vector space over } \mathbb{C}.
\]
\end{remark*}

\begin{remark*}[Definition predicate reading]
\[
\operatorname{ComplexVectorSpace}(V)
\;\colonequiv\;
\operatorname{VectorSpace}(V,\mathbb{C},+,\cdot).
\]
\end{remark*}

\begin{remark*}[Interpretation]
In a complex vector space, the scalars used in scalar multiplication are complex numbers.
\end{remark*}

\begin{dependencies}
\begin{itemize}
  \item \hyperref[def:vector-space]{Definition: Vector Space}
\end{itemize}
\end{dependencies}

\begin{remark*}[Notation]
For each positive integer $m$, the notation $\mathbf{F}^m$ denotes the vector
space of all lists of length $m$ of elements of $\mathbf{F}$.

The notation $\mathbf{F}^{\infty}$ denotes the set of all infinite lists
\[
(x_1,x_2,x_3,\dots)
\]
with each $x_j \in \mathbf{F}$, equipped with coordinatewise addition and
coordinatewise scalar multiplication.
\end{remark*}

\begin{remark*}[Coordinates in Vector Spaces]
A pointwise proof is available when vectors are given explicitly by coordinates, as in $\mathbf{F}^n$, because the operations are defined coordinatewise. In an arbitrary vector space, vectors need not come with coordinates, so proofs must use the vector space axioms rather than coordinate calculations.
\end{remark*}

% ---------------------------------------------------------

\begin{theorembox}{Theorem (Unique Additive Identity)}
\begin{theorem}[Unique Additive Identity]
\label{thm:unique-additive-identity}
A vector space has a unique additive identity.
\hyperref[prf:unique-additive-identity]{\textit{Go to proof.}}
\end{theorem}
\end{theorembox}

\begin{remark*}[Standard quantified statement]
\[
\begin{aligned}
&\text{If } 0,0' \in V \text{ each satisfy} \\
&\qquad \forall v \in V,\; v+0=v
\quad\text{and}\quad
\forall v \in V,\; v+0'=v, \\
&\text{then } 0=0'.
\end{aligned}
\]
\end{remark*}

\begin{remark*}[Interpretation]
Although the vector space axioms assert the existence of an additive identity,
they force that identity to be unique.
\end{remark*}

\begin{dependencies}
\begin{itemize}
  \item \hyperref[def:vector-space]{Definition: Vector Space}
\end{itemize}
\end{dependencies}

% ---------------------------------------------------------

\begin{theorembox}{Theorem (Unique Additive Inverse)}
\begin{theorem}[Unique Additive Inverse]
\label{thm:unique-additive-inverse}
Every vector in a vector space has a unique additive inverse.
\hyperref[prf:unique-additive-inverse]{\textit{Go to proof.}}
\end{theorem}
\end{theorembox}

\begin{remark*}[Standard quantified statement]
\[
\begin{aligned}
&\forall v \in V,\;
\text{if } w,u \in V \text{ satisfy } v+w=0 \text{ and } v+u=0, \\
&\quad\text{then } w=u.
\end{aligned}
\]
\end{remark*}

\begin{remark*}[Interpretation]
The vector space axioms guarantee that each vector has an additive inverse,
and that inverse is determined uniquely by the requirement that the sum be $0$.
\end{remark*}

\begin{remark*}[Notation]
If $v \in V$, then the unique additive inverse of $v$ is denoted by $-v$.

If $w,v \in V$, then
\[
w-v
\;:=\;
w+(-v).
\]
\end{remark*}

\begin{dependencies}
\begin{itemize}
  \item \hyperref[def:vector-space]{Definition: Vector Space}
  \item \hyperref[thm:unique-additive-identity]{Theorem: Unique Additive Identity}
\end{itemize}
\end{dependencies}

% ---------------------------------------------------------

\begin{theorembox}{Theorem ($0v=0$)}
\begin{theorem}[$0v=0$]
\label{thm:zero-scalar-times-vector}
For every $v \in V$,
\[
0v=0.
\]
\hyperref[prf:zero-scalar-times-vector]{\textit{Go to proof.}}
\end{theorem}
\end{theorembox}

\begin{remark*}[Standard quantified statement]
\[
\forall v \in V,\; 0v=0.
\]
\end{remark*}

\begin{remark*}[Interpretation]
Multiplying any vector by the scalar $0$ gives the zero vector. The $0$ on the
left is the scalar zero in $\mathbf{F}$; the $0$ on the right is the zero vector
in $V$.
\end{remark*}

\begin{dependencies}
\begin{itemize}
  \item \hyperref[def:vector-space]{Definition: Vector Space}
  \item \hyperref[thm:unique-additive-inverse]{Theorem: Unique Additive Inverse}
\end{itemize}
\end{dependencies}

% ---------------------------------------------------------

\begin{theorembox}{Theorem ($a0=0$)}
\begin{theorem}[$a0=0$]
\label{thm:scalar-times-zero-vector}
For every $a \in \mathbf{F}$,
\[
a0=0.
\]
\hyperref[prf:scalar-times-zero-vector]{\textit{Go to proof.}}
\end{theorem}
\end{theorembox}

\begin{remark*}[Standard quantified statement]
\[
\forall a \in \mathbf{F},\; a0=0.
\]
\end{remark*}

\begin{remark*}[Interpretation]
Every scalar sends the zero vector to the zero vector. The $0$ on the right of
$a$ is the zero vector in $V$.
\end{remark*}

\begin{dependencies}
\begin{itemize}
  \item \hyperref[def:vector-space]{Definition: Vector Space}
  \item \hyperref[thm:unique-additive-inverse]{Theorem: Unique Additive Inverse}
\end{itemize}
\end{dependencies}

% ---------------------------------------------------------

\begin{theorembox}{Theorem ($(-1)v=-v$)}
\begin{theorem}[$(-1)v=-v$]
\label{thm:minus-one-times-vector}
For every $v \in V$,
\[
(-1)v=-v.
\]
\hyperref[prf:minus-one-times-vector]{\textit{Go to proof.}}
\end{theorem}
\end{theorembox}

\begin{remark*}[Standard quantified statement]
\[
\forall v \in V,\; (-1)v=-v.
\]
\end{remark*}

\begin{remark*}[Interpretation]
Scalar multiplication by $-1$ produces the additive inverse of a vector.
This justifies the notation $-v$ for the additive inverse: it is literally
$(-1)$ times $v$.
\end{remark*}

\begin{dependencies}
\begin{itemize}
  \item \hyperref[def:vector-space]{Definition: Vector Space}
  \item \hyperref[thm:unique-additive-inverse]{Theorem: Unique Additive Inverse}
  \item \hyperref[thm:zero-scalar-times-vector]{Theorem: $0v=0$}
\end{itemize}
\end{dependencies}
